% =====================================================================
\section{Discussion}
\label{sec:discussion}
% =====================================================================

Complementarity here is a property of the task, not of imagery in general. An
attribute question---local land use, building height, transit access---is answerable
from one location's imagery, so the two views carry much the same answer-relevant
content and the combined input rarely beats the stronger one. A cross-view question
depends instead on the relationship between the views, where our models do use both.

The measurements show the model derives little usable signal from the second view,
and that the combined input is on net slightly worse than the better single view;
they do not identify the internal cause. Redundant information across views, a
representation dominated by one modality, and an architecture that fails to combine
them would all produce the pattern we observe, and our data do not distinguish among
them. The input--output fact---adding the second view does not help and often
hurts---is what a practitioner deciding whether to collect it needs to know.

This has a practical consequence. Street-level and high-resolution satellite imagery
differ sharply in cost, coverage, and licensing, yet for these attribute families the
second view adds at most a couple of points and degrades about one item in ten. A
deployment can serve a single view---the cheaper or better-covered one---at little
accuracy cost; and when both are available, the evidence favors \emph{routing} each
question to one view over feeding the model both.

\subsection{Limitations and scope}
\label{sec:limitations}
The claim concerns seven urban-attribute tasks in four-option MCQ form at commonly
served satellite resolutions; it is not a claim that overhead imagery is
uninformative or that fusion fails in general, as the cross-view tasks show large
synergy on the same models. Three caveats bound it further: our two training regimes
share one held-out test set, of which the held-out-cities split is the true
generalization check; higher ground sampling distance might help the few
overhead-geometry-limited tasks (building height, urban density); and an open-ended
or regression formulation could surface complementary signal an MCQ cannot express.

% =====================================================================
\section{Conclusion}
\label{sec:conclusion}
% =====================================================================

We measured, rather than assumed, the value of a second perspective for urban
attribute recognition. Across model scale, training regime, inference-time view
count, and prompting strategy, the second view adds at most a couple of points over
the better single view, and item by item it removes more correct answers than it
adds; the same models do benefit on tasks whose answer lives in the correspondence
between the two views. We release OELLM---benchmark, adapters, and protocol---for
reuse. The practical lesson is simple: before collecting a second visual modality,
ask whether the task structurally requires cross-view reasoning; if it does not,
current models will not turn the extra view into accuracy.

% =====================================================================
\section{Future Work}
\label{sec:future}
% =====================================================================

Because only the cross-view tasks consistently benefit from multiple street views,
the natural next step is to exploit those views for richer geometric reasoning. The
four road-aligned angles per location contain the signal needed for local 3D
reconstruction; lifting urban understanding from 2D image fusion to 3D spatial
reasoning would let a model reason about occlusion, scale, building geometry, and
pedestrian-environment affordances in ways single 2D images cannot support, with
OELLM's four-angle captures as a natural input substrate. Further directions include
fusion architectures that combine modalities more effectively on redundant-view
tasks, scaling to more cities to reduce the Europe-heavy distribution, and extending
deterministic OSM grounding to socio-economic indicators.
